\documentclass[11pt,a4paper]{article}

\usepackage[utf8]{inputenc}
\usepackage[T1]{fontenc}
\usepackage{geometry}
\geometry{a4paper, margin=2.5cm}
\usepackage{amsmath}
\usepackage{amssymb}
\usepackage{xcolor}
\usepackage[colorlinks=true, linkcolor=blue!50!black, urlcolor=blue!50!black, citecolor=blue!50!black]{hyperref}
\usepackage{enumitem}
\usepackage{booktabs}
\usepackage{titlesec}
\usepackage{fancyhdr}
\usepackage{graphicx}
\usepackage{caption}

% ---------- Heading style ----------
\titleformat{\section}{\large\bfseries}{\thesection}{0.6em}{}
\titleformat{\subsection}{\normalsize\bfseries}{\thesubsection}{0.6em}{}

% ---------- Header and footer ----------
\pagestyle{fancy}
\fancyhf{}
\renewcommand{\headrulewidth}{0.4pt}
\fancyhead[L]{Computer Games \textendash{} Phase B Project Plan}
\fancyhead[R]{Topic: Path Planning}
\fancyfoot[C]{\thepage}

\setlength{\parindent}{0pt}
\setlength{\parskip}{0.5em}

\title{\Huge\textbf{Project Plan}\\[0.3em]
\Large Computer Games \textendash{} Phase B\\[0.2em]
\large Topic: Path Planning}
\author{}
\date{}

\begin{document}

\maketitle
\thispagestyle{fancy}

\noindent
\textbf{Group:} \underline{1} \qquad
\textbf{Members:} Franziska Plagwitz, Luna, Sem Massa

\vspace{1em}

\section*{Overview}

This document defines our shared approach for the project in the Computer
Games Exercises, Phase B. Our assigned topic is \textbf{Path Planning}. The
goal is a small game or demo application that demonstrates at least three
techniques covered in the lecture and additionally includes one recently
developed technique. The focus lies on a direct, measurable comparison of
different path planning methods.

\section{Big Picture Goal}

We will build a \textbf{life simulation demo} in \textbf{Godot}: four NPCs
move through a top down house scene with several connected rooms, each
NPC following its own list of two to three targets it visits one after
another (looping once the list is finished). For every target, the NPC
runs through a small finite state machine, walking toward the target
(state \texttt{WALKING}), then interacting with it for a fixed duration
(state \texttt{IDLE}) before moving on. A switchable path planning
algorithm (default $A^*$) computes the route for every walk phase, so
the same scenario can be replayed with different algorithms and compared
directly. Benchmarks are recorded live during the simulation and
displayed as a comparison chart at the end.

\subsection{Game Concept in Detail}

\begin{itemize}[leftmargin=1.2em]
  \item \textbf{Scenario:} top down view of a house with multiple rooms
  connected by doors and narrow passages, populated with static furniture
  obstacles (see Figure~\ref{fig:house-layout}).
  \item \textbf{NPCs and targets:} four NPCs, each assigned an array of
  two to three target points inside the house. Each NPC cycles through
  \texttt{WALKING} (path planning and movement to the next target) and
  \texttt{IDLE} (a countdown timer at the target before moving on), and
  loops back to the start of its target array once finished.
  \item \textbf{Mode 1, Spectator Race:} all four NPCs are active
  simultaneously, each one possibly using a different path planning
  algorithm, so the resulting routes through the house can be compared
  visually as the simulation runs.
  \item \textbf{Mode 2, Live Benchmark:} on button press, a fixed set of
  start/target pairs is solved by several algorithms one after another,
  without NPC animation, purely for data collection.
  \item \textbf{Dynamic obstacles (optional extension):} some obstacles
  move on their own so that replanning becomes necessary. This makes the
  difference between algorithms visible and provides material for the
  recent technique.
\end{itemize}

\begin{figure}[h]
  \centering
  \includegraphics[width=0.85\textwidth]{Bildschirmfoto 2026-06-20 um 18.34.44.png}
  \caption{Sketch of the house scene with four NPCs (colored stars) and
  their target points (colored squares/circles), from the original
  project idea.}
  \label{fig:house-layout}
\end{figure}

\newpage 

\subsection{Proposed Techniques (Lecture Integration)}

The following points outline various techniques and lecture concepts that could be integrated into the project. These serve as suggestions and options for implementation rather than rigid specifications:

\begin{enumerate}[leftmargin=1.4em]
  \item \textbf{Path Planning \& Navigation (Core Topic):} 
  A direct comparison of algorithms (e.g., Dijkstra/BFS, Greedy Best-First, and $A^*$) could be conducted. For this, different discrete navigation abstractions mentioned in the lecture might be utilized, such as a uniform raster/grid, waypoint graphs, or polygon meshes (NavigationMesh).
  
  \item \textbf{Game Engine Architecture \& Messaging (Potential Approach):} 
  \begin{itemize}
    \item \textit{Messaging System:} To reduce system coupling instead of hardcoding direct dependencies, the live benchmark could utilize Godot's signal architecture as an asynchronous, event-driven message bus. Agents might emit detached data objects caught by a central \texttt{BenchmarkManager}.
    \item \textit{State Machines:} Finite state machines (FSM) represent a viable option to handle behavioral switches for the pursued actor and to track replanning phases.
  \end{itemize}
  
  \item \textbf{Collision Detection Phases \& Proxies (Optional Optimization):} 
  \begin{itemize}
    \item \textit{Broad Phase vs. Narrow Phase:} The optimization split could be explored by relying on the engine's broad-phase spatial partition hierarchies (Grid/BVH) to eliminate non-colliding pairs, while precise narrow-phase geometric tests handle near-field agent-to-obstacle resolution.
    \item \textit{Bounding Boxes \& Proxies:} Incorporating Axis-Aligned Bounding Boxes (AABB) or simple bounding circles as geometric proxies is another possibility. Dynamic obstacles changing their spatial occupancy could then trigger localized graph updates and instant path replanning.
  \end{itemize}
  
  \item \textbf{Graphics Pipeline \& Transformation Spaces (Suggested Concepts):} 
  \begin{itemize}
    \item \textit{Transform Spaces \& Model Matrices:} Agent coordinates could be mapped through distinct affine spaces. For instance, every agent might utilize localized 2D translation and rotation transformation matrices to project coordinate points into the shared Global/World Space.
    \item \textit{Scene Tree Hierarchy:} Godot's hierarchical tree node structure could be used to demonstrate parent-child relative transform propagation (e.g., obstacles attached to a moving compound platform inheriting spatial translation).
  \end{itemize}
  
  \item \textbf{Particle Simulations (Optional Visual Feature):} 
  \begin{itemize}
    \item To visibly differentiate competing path planning strategies, agents could optionally emit real-time, color-coded trailing paths utilizing particle emitters. This simulation might map to the lecture's four-step particle life cycle mechanics: 1) \textit{Produce/Inject} at current spatial nodes, 2) assigning individual \textit{Attributes} (initial reverse velocity, color tags), 3) scripted \textit{Transport/Modification}, and 4) deletion once time-to-live (\textit{TTL}) limits \textit{Vanish}.
  \end{itemize}
  
  \item \textbf{Key Frame Animation (Visual Polish):}
  Walking and idle states from the original idea could be backed by simple
  key frame animation cycles (e.g. a walk cycle as in the original
  sprite sheet), switched based on the NPC's current FSM state.
\end{enumerate}

\newpage

\section{Work Packages}

The three work packages are scoped so they can be tackled largely in
parallel. 

\subsection*{WP1, Game Foundation, Scene, and Collision Detection}

\textbf{Responsible:} \underline{\hspace{4cm}}

\textbf{Goal:} a playable Godot scene with obstacles, a navigation
structure, and working collision detection, on top of which the
algorithms from WP2 will operate.

\textbf{Tasks:}
\begin{itemize}[leftmargin=1.2em]
  \item set up the Godot project, build a scene with static obstacles
  \item choose and implement a navigation structure (grid, waypoint graph,
  or NavigationMesh, decided jointly with WP2 since this forms the basis
  for the algorithm comparison)
  \item collision detection: queries between agent and wall, and between
  agents themselves
  \item optional: dynamic obstacles with their own movement
  \item player character or flee logic for the target
  \item basic movement and steering of agents along a given path (path
  following)
\end{itemize}

\textbf{Interface to WP2:} provides the graph structure (nodes, edges,
weights) on which the pathfinding algorithms operate.

\subsection*{WP2, Path Planning Algorithms and Benchmarks}

\textbf{Responsible:} \underline{\hspace{4cm}}

\textbf{Goal:} implementation of the algorithms to be compared, together
with a benchmark framework that records metrics during the simulation.

\textbf{Tasks:}
\begin{itemize}[leftmargin=1.2em]
  \item literature research on recent papers for the recent technique, and
  selection of benchmark metrics (e.g. path length, number of expanded
  nodes, computation time, success rate under dynamic obstacles,
  replanning frequency)
  \item implement dynamic programming / breadth first search as baseline
  \item implement greedy best first search
  \item implement A* search
  \item implement the chosen recent technique (e.g. Jump Point Search;
  Particle Repulsion is an alternative candidate from the original idea)
  \item benchmark framework: logging of metrics per algorithm and run
  \item visualization of paths in the scene (e.g. color coded per
  algorithm)
\end{itemize}

\textbf{Interface to WP1:} requires the graph structure from WP1 as input,
returns computed paths back to agent steering.

\textbf{Interface to WP3:} provides the raw benchmark data for evaluation
and the report.

\subsection*{WP3, Evaluation, Report, and Presentation}

\textbf{Responsible:} \underline{\hspace{4cm}}

\textbf{Goal:} a clearly structured report with proper references that
explains the techniques and evaluates the benchmark results, plus the
final presentation.

\textbf{Tasks:}
\begin{itemize}[leftmargin=1.2em]
  \item structure of the report (outline, organization, LaTeX template)
  \item evaluation of the benchmark data from WP2: charts, tables,
  statistical assessment
  \item comparison with results from recent papers (assessing whether our
  results align with the literature)
  \item theory section: explanation of the techniques used, with
  references
  \item merging contributions from all three work packages into a
  consistent report
  \item preparation of the presentation (slides, live demo of the game)
  \item rehearsal presentation
\end{itemize}

\textbf{Note:} this work package starts somewhat later in time, since it
depends on initial results from WP1 and WP2, but runs in parallel
content wise (e.g. the theory section can already be started early).



\section{Open Decisions}

\begin{itemize}[leftmargin=1.2em]
  \item final choice of the recent technique (suggestion: Jump Point
  Search, since it compares well against A* and is well documented in
  recent literature)
  \item navigation structure: grid versus waypoint graph versus
  NavigationMesh (Godot provides NavigationMesh natively, which eases
  implementation but may offer less control over the graph structure for
  the algorithm comparison)
  \item scope of dynamic obstacles (purely static would be simpler,
  dynamic makes the comparison more interesting and demonstrates
  replanning)
\end{itemize}

\end{document}
